% GENERATED FILE. Edit canonical lessons, then run npm run generate:books.
% canonical-source-hash: fnv1a64:e116be913b2b232e
% canonical-lessons: KA-C343-nirupisu, KA-C343-varnisu, KA-C343-sannivesa, KA-C343-prasanga, KA-C343-hinnele

\chapter{How to Tell a Story}
\label{ch:ka-how-to-tell-a-story}
\hlchaptermodality{343}

\begin{chapteropening}
\textbf{By the end of this chapter:} \emph{I can say to narrate, to describe vividly, a situation, an episode and a background.}

Complete the last lesson of chapter 343: I can say to narrate, to describe vividly, a situation, an episode and a background.
\end{chapteropening}

\section[\texorpdfstring{\kn{ನಿರೂಪಿಸು} (nirūpisu)}{ನಿರೂಪಿಸು (nirūpisu)}]{\kn{ನಿರೂಪಿಸು} (nirūpisu) — to narrate; to present}

\label{lesson:KA-C343-nirupisu}



\begin{quote}
Before the new one: say the Kannada for to surrender, then the Kannada for a cannon.
\end{quote}

\subsection*{You'll want to know: \kn{ನಿರೂಪಿಸು}}
\textbf{\kn{ನಿರೂಪಿಸು}} — \emph{nirūpisu} — \textquotedblleft{}to narrate; to present\textquotedblright{}.

\begin{quote}
\kn{ಅಜ್ಜ} \kn{ಯುದ್ಧದ} \kn{ಕಥೆಯನ್ನು} \kn{ಚೆನ್ನಾಗಿ} \kn{ನಿರೂಪಿಸಿದರು}.
\end{quote}

\begin{quote}
\emph{ajja yuddhada katheyannu cennāgi nirūpisidaru.}
\end{quote}

\begin{quote}
Grandfather narrated the war story well.
\end{quote}

Past (I …): \kn{ನಿರೂಪಿಸಿದೆ}, \emph{nirūpiside}.

What is told takes -\kn{ಅನ್ನು}. A programme host is a \kn{ನಿರೂಪಕ}.

\begin{cousinweb}
\textbf{\kn{ನಿರೂಪಿಸು}} (\emph{nirūpisu}) is built on Sanskrit \emph{nirūpaṇa}, \textquotedblleft{}stating, defining\textquotedblright{} (Monier-Williams), with the Kannada verb ending -\emph{isu}.
\end{cousinweb}

\begin{grammarlens}[title={What you've built}]
One more word for telling what happened.
\end{grammarlens}

\subsection*{Your turn}
\begin{itemize}
\raggedright
  \item \emph{Say it:} \emph{nirūpisu}
  \item \emph{Say it:} \emph{nirūpisu}, once more
  \item \emph{Say it:} \emph{ajja yuddhada katheyannu cennāgi nirūpisidaru.} — Grandfather narrated the war story well.
  \item \emph{Recall:} say the Kannada for to surrender, then the Kannada for a cannon, then say \emph{nirūpisu} again
\end{itemize}

\begin{tcolorbox}[breakable,colback=teal!4,colframe=teal!35!black,title={Before you move on}]
What does \kn{ನಿರೂಪಿಸು} mean? (\textquotedblleft{}To narrate; to present\textquotedblright{}.) Say it once more.
\end{tcolorbox}
\section[\texorpdfstring{\kn{ವರ್ಣಿಸು} (varṇisu)}{ವರ್ಣಿಸು (varṇisu)}]{\kn{ವರ್ಣಿಸು} (varṇisu) — to describe vividly}

\label{lesson:KA-C343-varnisu}



\begin{quote}
Before the new one: say the Kannada for a cannon, then the Kannada for to narrate.
\end{quote}

\subsection*{You'll want to know: \kn{ವರ್ಣಿಸು}}
\textbf{\kn{ವರ್ಣಿಸು}} — \emph{varṇisu} — \textquotedblleft{}to describe vividly\textquotedblright{}.

\begin{quote}
\kn{ಅವಳು} \kn{ಹಳ್ಳಿಯ} \kn{ಜಾತ್ರೆಯನ್ನು} \kn{ಸುಂದರವಾಗಿ} \kn{ವರ್ಣಿಸಿದಳು}.
\end{quote}

\begin{quote}
\emph{avaḷu haḷḷiya jātreyannu sundaravāgi varṇisidaḷu.}
\end{quote}

\begin{quote}
She described the village fair beautifully.
\end{quote}

Past (I …): \kn{ವರ್ಣಿಸಿದೆ}, \emph{varṇiside}.

\kn{ವಿವರಿಸು} is to explain; \kn{ವರ್ಣಿಸು} is to paint in words.

\begin{cousinweb}
\textbf{\kn{ವರ್ಣಿಸು}} (\emph{varṇisu}) is built on Sanskrit \emph{varṇana}, \textquotedblleft{}description\textquotedblright{} (CDIAL 11340), from \emph{varṇa}, \textquotedblleft{}colour\textquotedblright{}, with the Kannada verb ending -\emph{isu}.
\end{cousinweb}

\begin{grammarlens}[title={What you've built}]
One more word for telling what happened.
\end{grammarlens}

\subsection*{Your turn}
\begin{itemize}
\raggedright
  \item \emph{Say it:} \emph{varṇisu}
  \item \emph{Say it:} \emph{varṇisu}, once more
  \item \emph{Say it:} \emph{avaḷu haḷḷiya jātreyannu sundaravāgi varṇisidaḷu.} — She described the village fair beautifully.
  \item \emph{Recall:} say the Kannada for a cannon, then the Kannada for to narrate, then say \emph{varṇisu} again
\end{itemize}

\begin{tcolorbox}[breakable,colback=teal!4,colframe=teal!35!black,title={Before you move on}]
What does \kn{ವರ್ಣಿಸು} mean? (\textquotedblleft{}To describe vividly\textquotedblright{}.) Say it once more.
\end{tcolorbox}
\section[\texorpdfstring{\kn{ಸನ್ನಿವೇಶ} (sannivēśa)}{ಸನ್ನಿವೇಶ (sannivēśa)}]{\kn{ಸನ್ನಿವೇಶ} (sannivēśa) — a situation; a scene}

\label{lesson:KA-C343-sannivesa}



\begin{quote}
Before the new one: say the Kannada for to narrate, then the Kannada for to describe vividly.
\end{quote}

\subsection*{You'll want to know: \kn{ಸನ್ನಿವೇಶ}}
\textbf{\kn{ಸನ್ನಿವೇಶ}} — \emph{sannivēśa} — \textquotedblleft{}a situation; a scene\textquotedblright{}.

\begin{quote}
\kn{ಆ} \kn{ಸಿನಿಮಾದಲ್ಲಿ} \kn{ಒಂದು} \kn{ತಮಾಷೆಯ} \kn{ಸನ್ನಿವೇಶ} \kn{ಇದೆ}.
\end{quote}

\begin{quote}
\emph{ā sinimādalli ondu tamāṣeya sannivēśa ide.}
\end{quote}

\begin{quote}
There is a funny scene in that film.
\end{quote}

Plural: \kn{ಸನ್ನಿವೇಶಗಳು}, \emph{sannivēśagaḷu}.

A scene in a story or a film, or a real situation: \kn{ಕಷ್ಟದ} \kn{ಸನ್ನಿವೇಶ}.

\begin{cousinweb}
\textbf{\kn{ಸನ್ನಿವೇಶ}} (\emph{sannivēśa}) is borrowed from Sanskrit \emph{saṃniveśa}, \textquotedblleft{}entering together; a seat, a position, a situation\textquotedblright{} (Monier-Williams).
\end{cousinweb}

\begin{grammarlens}[title={What you've built}]
One more word for telling what happened.
\end{grammarlens}

\subsection*{Your turn}
\begin{itemize}
\raggedright
  \item \emph{Say it:} \emph{sannivēśa}
  \item \emph{Say it:} \emph{sannivēśa}, once more
  \item \emph{Say it:} \emph{ā sinimādalli ondu tamāṣeya sannivēśa ide.} — There is a funny scene in that film.
  \item \emph{Recall:} say the Kannada for to narrate, then the Kannada for to describe vividly, then say \emph{sannivēśa} again
\end{itemize}

\begin{tcolorbox}[breakable,colback=teal!4,colframe=teal!35!black,title={Before you move on}]
What does \kn{ಸನ್ನಿವೇಶ} mean? (\textquotedblleft{}A situation; a scene\textquotedblright{}.) Say it once more.
\end{tcolorbox}
\section[\texorpdfstring{\kn{ಪ್ರಸಂಗ} (prasaṅga)}{ಪ್ರಸಂಗ (prasaṅga)}]{\kn{ಪ್ರಸಂಗ} (prasaṅga) — an episode; an incident}

\label{lesson:KA-C343-prasanga}



\begin{quote}
Before the new one: say the Kannada for to describe vividly, then the Kannada for a situation.
\end{quote}

\subsection*{You'll want to know: \kn{ಪ್ರಸಂಗ}}
\textbf{\kn{ಪ್ರಸಂಗ}} — \emph{prasaṅga} — \textquotedblleft{}an episode; an incident\textquotedblright{}.

\begin{quote}
\kn{ಅಜ್ಜಿ} \kn{ತಮ್ಮ} \kn{ಬಾಲ್ಯದ} \kn{ಒಂದು} \kn{ಪ್ರಸಂಗ} \kn{ಹೇಳಿದರು}.
\end{quote}

\begin{quote}
\emph{ajji tamma bālyada ondu prasaṅga hēḷidaru.}
\end{quote}

\begin{quote}
Grandmother told an episode from her childhood.
\end{quote}

Plural: \kn{ಪ್ರಸಂಗಗಳು}, \emph{prasaṅgagaḷu}.

A short happening worth telling, from life or from a story.

\begin{cousinweb}
\textbf{\kn{ಪ್ರಸಂಗ}} (\emph{prasaṅga}) is borrowed from Sanskrit \emph{prasaṅga}, \textquotedblleft{}attachment, connection, an occasion\textquotedblright{} (Monier-Williams).
\end{cousinweb}

\begin{grammarlens}[title={What you've built}]
One more word for telling what happened.
\end{grammarlens}

\subsection*{Your turn}
\begin{itemize}
\raggedright
  \item \emph{Say it:} \emph{prasaṅga}
  \item \emph{Say it:} \emph{prasaṅga}, once more
  \item \emph{Say it:} \emph{ajji tamma bālyada ondu prasaṅga hēḷidaru.} — Grandmother told an episode from her childhood.
  \item \emph{Recall:} say the Kannada for to describe vividly, then the Kannada for a situation, then say \emph{prasaṅga} again
\end{itemize}

\begin{tcolorbox}[breakable,colback=teal!4,colframe=teal!35!black,title={Before you move on}]
What does \kn{ಪ್ರಸಂಗ} mean? (\textquotedblleft{}An episode; an incident\textquotedblright{}.) Say it once more.
\end{tcolorbox}
\section[\texorpdfstring{\kn{ಹಿನ್ನೆಲೆ} (hinnele)}{ಹಿನ್ನೆಲೆ (hinnele)}]{\kn{ಹಿನ್ನೆಲೆ} (hinnele) — a background}

\label{lesson:KA-C343-hinnele}



\begin{quote}
Before the new one: say the Kannada for a situation, then the Kannada for an episode.
\end{quote}

\subsection*{You'll want to know: \kn{ಹಿನ್ನೆಲೆ}}
\textbf{\kn{ಹಿನ್ನೆಲೆ}} — \emph{hinnele} — \textquotedblleft{}a background\textquotedblright{}.

\begin{quote}
\kn{ಈ} \kn{ಕಥೆಯ} \kn{ಹಿನ್ನೆಲೆ} \kn{ಒಂದು} \kn{ಹಳೆಯ} \kn{ಯುದ್ಧ}.
\end{quote}

\begin{quote}
\emph{ī katheya hinnele ondu haḷeya yuddha.}
\end{quote}

\begin{quote}
This story's background is an old war.
\end{quote}

\kn{ಹಿನ್}, behind, and \kn{ನೆಲೆ}, ground: what lies behind. \kn{ಹಿನ್ನೆಲೆ} \kn{ಸಂಗೀತ} is background music.

\begin{cousinweb}
\textbf{\kn{ಹಿನ್ನೆಲೆ}} (\emph{hinnele}) is a native Dravidian compound: \emph{hin}, \textquotedblleft{}behind\textquotedblright{} (DEDR 4205), and \emph{nele}, \textquotedblleft{}a footing, a place\textquotedblright{} (DEDR 3675).
\end{cousinweb}

\begin{grammarlens}[title={What you've built}]
One more word for telling what happened.
\end{grammarlens}

\subsection*{Your turn}
\begin{itemize}
\raggedright
  \item \emph{Say it:} \emph{hinnele}
  \item \emph{Say it:} \emph{hinnele}, once more
  \item \emph{Say it:} \emph{ī katheya hinnele ondu haḷeya yuddha.} — This story's background is an old war.
  \item \emph{Recall:} say the Kannada for a situation, then the Kannada for an episode, then say \emph{hinnele} again
\end{itemize}

\begin{tcolorbox}[breakable,colback=teal!4,colframe=teal!35!black,title={Before you move on}]
What does \kn{ಹಿನ್ನೆಲೆ} mean? (\textquotedblleft{}A background\textquotedblright{}.) Say it once more.
\end{tcolorbox}
